\documentclass[12pt,a4paper]{article}
\usepackage{amsmath,amsthm,amssymb}
\usepackage[utf8]{inputenc}
\usepackage[russian]{babel}

\makeatletter

\textheight=251mm %+20mm
\textwidth=163.4mm
\oddsidemargin=-2.1mm
\evensidemargin=-2.1mm
\topmargin=-16.5mm %-20mm

\newcommand*{\ralph}[1]{\expandafter\rralph\csname c@#1\endcsname}
\newcommand*{\rralph}[1]%
{\ifcase #1\or а\or б\or в\or г\or д\or е\or ж\or з\or и\or
к\or л\or м\or н\or о\or п\or р\or с\or т\or у\or
ф\or х\or ц\or ч\or ш\or щ\or э\or ю\or
я\fi}

\newcounter{problem}
\setcounter{problem}{4}
\newcounter{item}
\newcommand*{\z}{\par\medskip\addtocounter{problem}{1}\setcounter{item}{0}
\textbf{Задача \arabic{problem}.\ }}
\newcommand*{\zh}{\par\medskip\addtocounter{problem}{1}\setcounter{item}{0}
\textbf{Задача \arabic{problem}*.\ }}
\newcommand*{\zH}{\par\medskip\addtocounter{problem}{1}\setcounter{item}{0}
\textbf{Задача \arabic{problem}**.\ }}
\newcommand*{\zn}[1]{\par\medskip\addtocounter{problem}{1}\setcounter{item}{0}
\textbf{Задача \arabic{problem} (#1).}}
\newcommand*{\znh}[1]{\par\medskip\addtocounter{problem}{1}\setcounter{item}{0}
\textbf{Задача \arabic{problem}* (#1).}}
\newcommand*{\znH}[1]{\par\medskip\addtocounter{problem}{1}\setcounter{item}{0}
\textbf{Задача \arabic{problem}** (#1).}}

\newcommand*{\zite}{\par\medskip\addtocounter{problem}{1}\setcounter{item}{1}
\textbf{Задача \arabic{problem}.} (\ralph{item})\ }
\newcommand*{\ite}{\par\addtocounter{item}{1}(\ralph{item})\ }
\newcommand*{\iteh}{\par\addtocounter{item}{1}(\ralph{item})*\ }
\newcommand*{\iteH}{\par\addtocounter{item}{1}(\ralph{item})**\ }

\newcommand*{\iten}[1]{\par\addtocounter{item}{1}(\ralph{item}) \textbf{(#1)}}
\newcommand*{\itenh}[1]{\par\addtocounter{item}{1}(\ralph{item})* \textbf{(#1)}}
\newcommand*{\itenH}[1]{\par\addtocounter{item}{1}(\ralph{item})** \textbf{(#1)}}
\newcommand{\myP}[0]{\mathsf{P}}

\begin{document}
\pagestyle{empty}

%\renewcommand{\@oddhead}%
%{\vbox{
{\footnotesize\begin{tabbing}V зимняя школа <<Комбинаторная математика и теория алгоритмов>>\`10--17 февраля 2013 г.\\
Виталий Павленко, ФИВТ МФТИ\`\texttt{vk.com/vitalypavlenko}\end{tabbing}}

%}}

\begin{center}
\vskip 24pt
{\bf\LargeТеорема Байеса -- 2: Акинатор}
\end{center}

Для получения зачёта нужно решить предыдущий листочек. Кроме того, нужно либо решить этот листочек, либо написать программу, которая выигрывает в <<Быки и коровы>>, либо совершить оба этих поступка.

\z Пусть имеется источник, который выдаёт символы $a$, $b$ и $c$ с вероятностями $\myP(a) = 0.6$, $\myP(b) = 0.3$ и $\myP(c) = 0.1$.
\ite Какова энтропия этого источника? Что говорит нам знание энтропии о потенциальной возможности сжатия длинных сообщений источника?
\ite Предложите какой-нибудь простой префиксный код для этого источника, который не использует группировку символов в пары. Пусть кодируется исходное сообщение длины $N$. Какова ожидаемая длина закодированного сообщения?
\ite Попробуйте предложить какой-нибудь более сложный префиксный код, который использует группировку символов в пары. Удалось ли вам уменьшить ожидаемую длину закодированного сообщения?

\z Пусть мы играем в <<Быки и коровы>>. Покажите, как можно вычислить условную энтропию 
$$ H(\text{число} \mid (Q_1 = 1234, A_1 = \text{1 б., 2 к.}), (Q_2 = 5678, A_2 = \text{1 б., 0 к.}), (Q_3 = 3428, A_3 = \text{?})). $$

\z Пусть мы верим в закон Хика и ставим эксперимент, в котором человек повторяет за компьютером числа на скорость. Сначала компьютер показывал человеку длинную серию, в которой всего два различных числа, и среднее время реакции человека составило 0.2 секунды. Потом компьютер показывал человеку серию, в которой четыре различных числа, и среднее время реакции человека составило 0.35 секунды. Компьютер готовится показывать человеку серию из шести различных чисел. Какое среднее время реакции следует ожидать?

\z Рассмотрим два источника, каждый из которых порождает два различных символа. Источник $p$~--- это честная монетка: $p_1 = 0.5$, $p_2 = 0.5$. Источник $q$~--- это нечестная монетка: $q_1 = 0.7$, $q_2 = 0.3$. Что неэффективнее: кодировать честный источник кодом, оптимальным для нечестного источника, или наоборот?

\end{document}